\Needspace{18\baselineskip}
\section{Paths, renaming, and isomorphism}
\Needspace{24\baselineskip}
\subsection{Paths and their behaviour under plugging}
\begin{definition}[Adjacency and path]\label{term:path}
Edges $e,f$ are adjacent when some vertex belongs to both $t(e)$ and $s(f)$.
A path is a nonempty edge list: a singleton is a path; prepending $e$ to a path
starting at $f$ gives a path when $e,f$ are adjacent. Edges may repeat, and this
notion does not specify a chosen boundary vertex at either end.
\end{definition}
\lean{Graph/Paths}{11}{12}{Net.Adjacent}
\lean{Graph/Paths}{15}{17}{Net.Path}
For a list beginning $e,f$, the path predicate is equivalent to adjacency of
$e,f$ together with the path predicate for the tail starting at $f$. Invert the
path constructor for the forward implication and apply it for the reverse one.
\lean{Graph/Paths}{19}{23}{Net.path_pair_iff}
Any map on edges preserving adjacency maps paths to paths. Induct on the path:
singletons remain singletons, and the inductive step transports its first
adjacency and then its tail path. If the edge map preserves and reflects
adjacency, it preserves and reflects paths. Induct on the list: the empty list
is a path on neither side, singletons are paths on both sides, and longer lists
reduce to the first adjacency equivalence and the tail induction hypothesis.
\lean{Graph/Paths}{25}{31}{Net.Path.map}
\lean{Graph/Paths}{33}{42}{Net.path_map_iff}

\begin{lemma}[Host and inserted adjacencies]
Two surviving host edges are adjacent after plugging exactly when they were
adjacent before plugging. Two inserted edges that were adjacent remain adjacent.
\end{lemma}
\begin{proof}
For host reflection, pull both incidence witnesses back through their mapped
lists. Equality of their host images implies equality of the original vertices
by injectivity of $j$, giving an old shared witness. For host preservation map
an old shared witness through $j$. For inserted preservation map the shared
replacement witness through $k$; injectivity is unnecessary in this direction.
\end{proof}
\lean{Graph/Paths}{51}{61}{Net.Replacement.adjacent_host_iff}
\lean{Graph/Paths}{63}{66}{Net.Replacement.adjacent_insert}
\begin{definition}[Non-self-loop hole]\label{term:non-self-loop}
The hole $h$ is non-self-looping if no vertex belongs to both $s(h)$ and $t(h)$.
This requires the two lists to be disjoint; it is stronger than distinctness
within each individual list, already required of every net.
\end{definition}
\lean{Graph/Paths}{70}{70}{Net.Replacement.NonSelfLoop}
For such a hole, $k$ is globally injective. Equal internal images reduce to equal
names; mixed internal/boundary cases contradict freshness. Two input images or
two output images reduce to the appropriate port injection. Equality of an
input image and an output image would exhibit a host vertex in both $s(h)$ and
$t(h)$, contrary to non-self-looping.
\lean{Graph/Paths}{72}{93}{Net.Replacement.insert_injective}
Consequently inserted adjacency is reflected as well as preserved: pull the two
mapped-list witnesses back and use global injectivity of $k$ to identify them.
The other implication is the unconditional preservation result above.
\lean{Graph/Paths}{95}{104}{Net.Replacement.adjacent_insert_iff}

\begin{theorem}[Path preservation and reflection]
A list of surviving host edges is a path after plugging if and only if its old
edge list is a path. Any replacement path gives an inserted path. For a
non-self-loop hole the converse holds for inserted paths too.
\end{theorem}
\begin{proof}
For host paths, induct on the surviving-edge list. Empty and singleton lists
are immediate; a longer list reduces to host adjacency equivalence and the tail
hypothesis. For inserted preservation, apply path mapping to inserted adjacency
preservation. For inserted reflection under the stated restriction, apply the
general path equivalence lemma to inserted adjacency equivalence.
\end{proof}
\lean{Graph/Paths}{109}{120}{Net.Replacement.path_host_iff}
\lean{Graph/Paths}{122}{125}{Net.Replacement.path_insert}
\lean{Graph/Paths}{129}{132}{Net.Replacement.path_insert_iff}

\Needspace{24\baselineskip}
\subsection{Transport along carrier bijections}
A bijection on internal names extends to ports by mapping only internal vertices.
Its inverse laws follow by cases on a port: internal cases use the original
inverse laws and boundary cases are reflexive. Mapping the identity is pointwise
the identity, also by port cases; function extensionality gives the corresponding
identity of functions.
\lean{Graph/Isomorphism}{9}{14}{Port.bijection}
\lean{Graph/Isomorphism}{16}{17}{Port.map_id}
\lean{Graph/Isomorphism}{19}{20}{Port.map_id_fun}
Fix bijections $a:G_V\cong V'$ and $b:G_E\cong E'$. Reindexed sources and targets
at $e'\in E'$ are the old incidence lists at $b^{-1}(e')$, mapped by the port
extension of $a$.
\lean{Graph/Isomorphism}{35}{36}{Net.Reindex.source}
\lean{Graph/Isomorphism}{38}{39}{Net.Reindex.target}
Source uniqueness survives this transport: two shared mapped witnesses are equal
by port injectivity, so old source uniqueness identifies their old edge indices;
apply $b$ and cancel $b\circ b^{-1}$. The target uniqueness proof uses old target
uniqueness and the same cancellation.
\lean{Graph/Isomorphism}{41}{48}{Net.Reindex.source_unique}
\lean{Graph/Isomorphism}{50}{57}{Net.Reindex.target_unique}
For source coverage, a new source witness pulls back to an old source vertex,
which is not an output; port maps preserve this property. Conversely, pull a
non-output vertex backwards through $a^{-1}$, choose an old source edge by
coverage, and map both the edge and vertex forwards. The two inverse laws prove
that this is a witness for the original new vertex. Target coverage repeats this
argument with non-input vertices and old target coverage.
\lean{Graph/Isomorphism}{59}{70}{Net.Reindex.source_coverage}
\lean{Graph/Isomorphism}{72}{83}{Net.Reindex.target_coverage}
Give an internal name $w$ the label of $a^{-1}(w)$. Labelling an old vertex after
mapping by $a$ then agrees with its old label: the internal case cancels
$a^{-1}\circ a$, and boundary labels are fixed.
\lean{Graph/Isomorphism}{85}{88}{Net.Reindex.label_map}
\begin{definition}[Reindexed graph]\label{term:reindexing}
The reindexed net uses the new carriers, inverse-transported labels, and the
incidence lists just defined.
\end{definition}
\begin{proof}[Well-formedness]
Transport finite enumerations along the bijections. Distinctness follows from
port injectivity, uniqueness and coverage from the preceding lemmas. For typing,
compose incidence mapping with the vertex-label map, use label commutation, and
apply the original typing equation. The boundary lists are unchanged.
\end{proof}
\lean{Graph/Isomorphism}{94}{121}{Net.reindex}

\Needspace{24\baselineskip}
\subsection{Graph isomorphism and shallow context congruence}
\begin{definition}[Isomorphism]\label{term:isomorphism}
An isomorphism $G\cong H$ between nets with the same signature and boundary lists
is a bijection of internal vertices and a bijection of edges preserving internal
labels, edge labels, and both ordered incidence lists. Its port action fixes
boundary positions.
\end{definition}
\lean{Graph/Isomorphism}{123}{129}{Net.Iso}
Identity bijections and identity list maps yield a reflexive isomorphism. For the
inverse isomorphism, invert both bijections. Inverse label preservation follows
by applying the original equation at an inverse image and cancelling. To invert
an incidence equation, map its two sides by the inverse port map and cancel the
forward/inverse maps pointwise. For composition, compose carrier bijections;
compose label equalities, and use successive list mapping followed by the two
incidence-preservation equations.
\lean{Graph/Isomorphism}{135}{141}{Net.Iso.refl}
\lean{Graph/Isomorphism}{143}{157}{Net.Iso.symm}
\lean{Graph/Isomorphism}{159}{175}{Net.Iso.trans}
Isomorphisms preserve and reflect adjacency: transport a shared witness forwards
for preservation; for reflection pull both witnesses back from the incidence
maps and identify their preimages by port injectivity. Applying the general path
equivalence lemma then preserves and reflects paths, with their edge lists
mapped by the edge bijection.
\lean{Graph/Isomorphism}{177}{191}{Net.Iso.adjacent_iff}
\lean{Graph/Isomorphism}{193}{195}{Net.Iso.path_iff}
The original graph is isomorphic to its reindexing. Use the chosen carrier
bijections; inverse cancellation proves label preservation and each incidence
identity directly from the definition of reindexing.
\lean{Graph/Isomorphism}{199}{207}{Net.reindexIso}

Suppose $f:H\cong K$ and both graphs have the interface of host hole $h$.
The intended isomorphism between $C[h:=H]$ and $C[h:=K]$ is the identity on
host internal names and $f$ on replacement internal names, and likewise the
identity on surviving host edges and $f$ on inserted edges. Sum bijections
supply their inverse laws.
\lean{Graph/PlugIsomorphism}{16}{17}{Net.PlugIso.vertices}
\lean{Graph/PlugIsomorphism}{19}{20}{Net.PlugIso.edges}
The resulting port map commutes with host embedding: check the three port
constructors, all of which are reflexive. It commutes with insertion too:
internal vertices use the replacement part of the sum map, and boundary vertices
reduce to the host-embedding equation at the corresponding hole port.
\lean{Graph/PlugIsomorphism}{22}{24}{Net.PlugIso.map_host}
\lean{Graph/PlugIsomorphism}{26}{32}{Net.PlugIso.map_insert}
\begin{theorem}[Congruence of shallow plugging]
For every typed host and hole, $H\cong K$ implies $C[h:=H]\cong C[h:=K]$.
No non-self-loop hypothesis is required.
\end{theorem}
\begin{proof}
Use the sum bijections above. For each label equation split the carrier tag:
host labels are unchanged and inserted labels use $f$'s corresponding equation.
For a host source list, commute the new port map past host embedding pointwise.
For an inserted source list, commute it past insertion and then use source
preservation by $f$. Repeat with target lists and target preservation. These
arguments prove all isomorphism fields without reflecting through insertion.
\end{proof}
\lean{Graph/PlugIsomorphism}{38}{85}{Net.plugIso}
Here edge labels are preserved by literal equality. For boxed labels containing
bodies, this is not yet the recursive notion that would identify independently
renamed but isomorphic bodies. No such stronger congruence is claimed.
